\chapter{Not important}
\section{regression models}
\begin{definition}[Mean Regression model]
	The mean regression model describes the dependance of Y on X.
	\begin{equation}
		Y = f(X) + \epsilon
	\end{equation}
	The equation is called the regression equation. In the mean regression the expected value of error is 0.
\end{definition}

\subsubsection*{Ingredients}
\begin{enumerate}
	\item \textbf{Observations}: $Y = (Y_1, \dots, Y_n)$. Data is a collection of pairs of $(Y_i, X_i)$. $Y_i$s are called dependant or explained variables. $X_i$s are called yvb or independent variable or features. 
		\begin{itemize}
			\item Example 1: $X$ collection of measurements and $Y_i \in \{0, 1\}$ for patients. 
			\item Example 2: $Y_i$ is a value of process of time $t_i$. Like stock price.
		\end{itemize}  $X$ is called \textbf{design}, it could be either random or deterministic. We focus on deterministic design but the random design is more typical, for example in machine learning. 
	\item \textbf{Errors}: $\epsilon_i = Y_i - f(X_i)$. Usually(mean regression) $\mathbb{E} \epsilon_i = 0$. Denote $\sigma_i^2 = Var[\epsilon_i]$. We focus on the homogeneous case(the variance does not dependant on i). Most strict case: homogenous Gaussian noise.
	\item \textbf{Regression function}: f(x). Examples: 
		\begin{itemize}
			\item linear
			\item polynomial
			\item linear basis expansion for a collection of functions $\{\psi_i \}_i$
			\item parametric regression
				\begin{equation*}
					f(x) \in \{ f(x, \theta), \theta \in \Theta \subset \mathbb{R}^p \}
				\end{equation*}
				$\exists \theta^*$ s.t.
				\begin{equation}
					f(x) = f(x, \theta^*), \quad Y_i = f(X_i, \theta^*) + \epsilon_i
				\end{equation} 
			\end{itemize}
\end{enumerate}




In the regression model, our goal is to retract information in $x$ about $y$. In the mean regressions we have centralised error($\mathbb{E} Y_i = f(x_i)$), $f$ is the regression function(systematic), allows to predict and forecast. So our aim is to recover $f$.\par 
But in reality, we always have to deal with the difference between real distribution and regression model. We have 2 main concerns.1)Distribution error; 2)Form of the regression function.\par 
We have the model: $Y_i = f(x_i) + \epsilon_i$, where $f$ is the systematic part and $\epsilon$ the individual error. Usually we assume(mean regression): $\mathbb{E} \epsilon_i = 0$ to be true. Later we make our study under the assumption $\epsilon_i$ is Gaussian. In the analysis we try to account for the real data distribution. Model is used to derive the methods but the analysis is under true distribution.\par 
As for the regression function $\mathbb{E} Y_i = f(X_i)$. The typical assumptions including:
\begin{itemize}
	\item $f(x) = f(x, \theta^*)$
	\item $f(x) = \sum_{i = 1}^p \theta_j \psi_j(x), \quad p \leq \infty$
\end{itemize}

Our approach:
\begin{itemize}
	\item design some methods for the parametric model $Y_i = f(x_i) + \epsilon_i$;
	\item Study and analysis under $\mathbb{P}$.
\end{itemize} 




